\section{Institutional Motivation: Procurement as Industrial Policy}
\label{sec:institutional}

Public procurement is not only a purchasing technology. It can also determine which firms, technologies, and local production capabilities are sustained by public demand. The industrial-policy role of procurement is particularly visible in innovation-oriented procurement, local-content rules, strategic supply-chain initiatives, and government efforts to build domestic or regional production capacity. The recent survey by \citet{ChiappinelliGiuffridaSpagnolo2025} treats public procurement explicitly as an innovation-policy instrument and emphasizes design, frictions, and its interaction with the broader policy mix. The OECD likewise documents procurement measures used to support domestic and local production, employment, strategic sectors, and supply-chain resilience \citep{OECD2026}.

A recurrent policy tension is that competition and localization need not have identical incidence. In the Japanese local-government setting, \citet{SaitoMitsuta2012} examine regional eligibility requirements that promote local industry and employment while restricting outside participation. Competitive tendering can reduce expected procurement costs and select more efficient suppliers, yet it can also compress supplier rents. Procurement rules that favor local or small suppliers make this tension explicit. \citet{Vagstad1995} analyzes procurement discrimination when governments care about local firms' profits. \citet{Marion2007} and \citet{AtheyCoeyLevin2013} show empirically and structurally that procurement preferences and set-asides change entry, bidder profit, procurement cost, and efficiency. These papers study procurement design. Our paper instead holds the procurement rule fixed and asks how the competitiveness of the supplier market feeds back into a separate government decision over industrial-policy roles.

The distinction is useful for regional policy. Suppose a cross-regional project has a lead or buyer side, denoted $U$, and a complementary supplier-side industrial base, denoted $D$. A differentiated policy profile means that one jurisdiction supports the lead/buyer role while the other hosts or develops the supplier base from which procurement is conducted. If the supplier-side regional government internalizes only part of the supplier surplus generated by that procurement relationship, then the incidence of procurement surplus affects its willingness to specialize. A highly competitive procurement market may be desirable for a central buyer while reducing the local payoff that sustains decentralized specialization.

The paper therefore does not model procurement competition as a government choice. The number of suppliers $N$ is an exogenous market-structure parameter. This keeps the object distinct from the bidder-attraction problem studied by \citet{Szech2011} and \citet{XuLi2019}, and from endogenous-entry procurement models such as \citet{LiZheng2009}. The comparative static asks how different supplier-market structures change the equilibrium of the regional policy game. Endogenizing supplier entry would introduce an additional participation equilibrium and is outside the model.

This focus also matches the recent industrial-organization framing of industrial policy. \citet{VanBiesebroeckZhang2026} emphasize that industrial-policy effects depend on prevailing market structure, the localization of supply chains, and coordination across market actors. \citet{MottaPolo2026} likewise show that the appropriate industrial-policy response to supply-chain risk depends on market organization and appropriability. The present paper isolates a different channel: market competition affects the geographic incidence of rents, and that incidence changes decentralized policy specialization.
